\section{Testbed and Experimental Design}
\label{sec:system}

The testbed is the \mosaic architecture of~\cite{mosaic2026}: two
transport stages, depicted in Figure~\ref{fig:architecture} and
summarised in Table~\ref{tab:testbed}. Stage~1 delivers encoded \vpcc
content from the origin server to an edge renderer over \quic, and is
where controlled impairment is injected. Stage~2 casts the rendered
stereo frames from the edge renderer to the user's \hmd over a private
\fiveg standalone network and is held constant across all conditions,
so that any variation in perceived quality originates in the Stage~1
transport path.

\begin{figure*}[t]
  \centering
  \resizebox{\textwidth}{!}{%
  \begin{tikzpicture}[
      font=\footnotesize,
      box/.style={draw, rounded corners=2pt, align=center, minimum height=10mm, inner sep=3pt, fill=white},
      stage/.style={draw, dashed, rounded corners=4pt, inner sep=5pt},
      slice/.style={draw, rounded corners=1pt, inner sep=2pt, font=\scriptsize, fill=white, text width=30mm, align=left},
      arr/.style={-{Latex[length=2mm]}, thick},
      lbl/.style={font=\scriptsize, text=black!70, align=center}
    ]
    \node[box, text width=26mm] (origin) {\textbf{Volumetric server}\\Caddy, \glsentryshort{http3}/\quic\\per-\glsentryshort{gof} chunks};
    \node[box, text width=26mm, right=18mm of origin] (tc) {\textbf{Impairment}\\Linux \glsentryshort{tc}\\(TBF + netem)};
    \node[box, text width=28mm, right=18mm of tc] (edge) {\textbf{Edge renderer}\\Unity~6, V3C decoder\\per-\glsentryshort{gof} telemetry};
    \node[box, text width=24mm, right=18mm of edge] (gnb) {\textbf{Private 5G SA}\\Ericsson\\RU / BBU / NC};
    \node[box, text width=24mm, right=14mm of gnb] (hmd) {\textbf{Meta Quest~3}\\5G USB dongle\\\glsentryshort{acr} rating};
    \draw[arr] (origin) -- node[above, lbl] {\quic v1\\one port per slice} (tc);
    \draw[arr] (tc) -- node[above, lbl] {\glsentryshort{http3} GET\\per \glsentryshort{gof}} (edge);
    \draw[arr] (edge) -- node[above, lbl] {ALVR\\stereo frames} (gnb);
    \draw[arr] (gnb) -- node[above, lbl] {5G NR} (hmd);
    \node[slice, below=4mm of tc] (sA) {\textbf{A} clean \hfill RA, \glsentryshort{gof}\,32};
    \node[slice, below=1mm of sA] (sB) {\textbf{B} loss \hfill \glsentryshort{ldp}, \glsentryshort{gof}\,4};
    \node[slice, below=1mm of sB] (sC) {\textbf{C} latency \hfill \glsentryshort{ldp}, \glsentryshort{gof}\,4};
    \node[slice, below=1mm of sC] (sD) {\textbf{D} loss + latency \hfill \glsentryshort{ldp}, \glsentryshort{gof}\,4};
    \draw[-, thin, black!60] (tc.south) -- (sA.north);
    \begin{scope}[on background layer]
      \node[stage, fit=(origin)(tc)(edge)(sD), label={[lbl]below:{Stage 1: transport path under test}}] {};
      \node[stage, fit=(gnb)(hmd), label={[lbl]below:{Stage 2: held constant}}] {};
    \end{scope}
  \end{tikzpicture}}
  \caption{End-to-end testbed. Stage~1 delivers \vpcc content from the
  origin to the edge renderer over reliable \quic through one of four
  impairment slices, each bound to its own UDP port; Stage~2 casts the
  rendered stereo frames to the headset over the private \fiveg network
  and is identical across conditions. All network effects of interest
  are observable at the edge renderer.}
  \label{fig:architecture}
\end{figure*}

\subsection{Stage 1: \vpcc Delivery over Reliable \quic}
\label{sec:sys-stage1}
The \vpcc-encoded sequences are hosted on a volumetric server running
Caddy~v2 (quic-go), which serves per-\gof chunks over \gls{http3} with
\gls{tls}~1.3. Chunks traverse a Linux Traffic Control chain (TBF + netem)
that implements the impairment profile of each slice; each slice is
bound to a distinct \gls{udp} port on the same origin, so all slices are
served from one Caddy instance without cross-condition interference.
The edge renderer is a Unity~6 application with a paced
parallel-download controller (up to three \gofs in flight,
\SI{200}{\milli\second} prefetch margin, per-\gof wall-clock
deadline). Arriving chunks are fed to the optimised native \gls{v3c} decoder
of the 5G-MAG platform~\cite{fivegmag}, integrated through a native
plugin that bypasses Unity's OpenXR path, and the reconstructed point
cloud is rendered locally. A \gof that misses its deadline is
substituted with the last good \gof, which is the mechanism that
produces freezes and frame skips. Per-\gof telemetry (request time,
completion time, deadline outcome) is logged for the Layer~1 analysis.
Loss accounting is exact at the IP layer; no partial-reliability or
datagram extension is used, so the transport behaviour is that of
standard \quic~\cite{rfc9000}.

\subsection{Stage 2: Edge-to-\hmd Delivery over Private \fiveg}
\label{sec:sys-stage2}
The edge renderer is colocated with the gNodeB of an Ericsson private
\fiveg standalone network. Rendered stereo frames are cast to the Meta
Quest~3 with ALVR through an OpenGL$\to$DX11 zero-copy interop pipeline
routed via SteamVR. Because the Quest~3 has no native 5G radio, a
commercial 5G USB dongle is adapted to its USB-C port and exposed to
the headset as a network interface, so the last mile traverses the 5G
radio access network rather than Wi-Fi~\cite{vquest2026}. Throughput
on this link is characterised with on-device capture (PCAPdroid) and
is not a bottleneck for the ALVR stream.

\begin{table}[t]
  \caption{Testbed components.}
  \label{tab:testbed}
  \centering
  \small
  \begin{tabular}{@{}ll@{}}
    \toprule
    Component & Specification \\
    \midrule
    Volumetric server & Ubuntu 24.04, Caddy v2.11 (quic-go), HTTP/3, TLS~1.3 AES-128-GCM \\
    Impairment        & Linux TC (TBF + netem chain), one UDP port per slice \\
    Edge renderer     & Windows~11, NVIDIA RTX~4070, Unity~6, 5G-MAG V3C native decoder \\
    VR casting        & ALVR, OpenGL$\to$DX11 interop via SteamVR \\
    5G network        & Ericsson private 5G SA, colocated edge \\
    \hmd              & Meta Quest~3 (Snapdragon XR2 Gen~2), adapted 5G USB dongle \\
    \bottomrule
  \end{tabular}
\end{table}

\subsection{Content and Encoder Configuration}
\label{sec:sys-content}
Three sequences from the 8i Voxelized Full Bodies dataset~\cite{8i},
Category~2 of the \vpcc common test conditions~\cite{vpccctc}, are
used: Longdress (vox10, $\approx$0.85\,M points/frame), Soldier
(vox10, $\approx$1.08\,M) and ThaiDancer (vox12, $\approx$3.13\,M), each
300 frames at \SI{30}{fps}. Sequences are encoded with \gls{tmc2}
v25~\cite{tmc2} using x265 for the geometry and attribute videos and
\gls{hm} for the losslessly coded occupancy maps (Table~\ref{tab:encoder}).
Attribute QPs are chosen per sequence to place all three in the same
\gls{psnr}$_Y$ band (\SIrange{34}{39}{dB}), so that differences across
conditions are attributable to transport rather than compression;
\pcqm is reported per sequence as a stimulus descriptor.

\begin{table}[t]
  \caption{\vpcc encoder configuration.}
  \label{tab:encoder}
  \centering
  \small
  \begin{tabular}{@{}ll@{}}
    \toprule
    Parameter & Value \\
    \midrule
    Codec                          & MPEG TMC2 v25.0 \\
    HEVC (geometry / attribute)    & x265 3.5 via HM-compatible shim \\
    HEVC (occupancy)               & HM 16.20 + SCM 8.8 \\
    Geometry QP / occupancy precision & 16 / 2 (fixed) \\
    Attribute QP (Soldier / Longdress / ThaiDancer) & 18 / 22 / 32 \\
    Clean-reference profile        & \glsentryshort{ra}, \gof{}=32, I/P/B frames \\
    Factorial profile              & LDP, \gof{}=4, I/P frames \\
    Structural contrast            & LDP, \gof{}=16, I/P frames \\
    \bottomrule
  \end{tabular}
\end{table}

\subsection{Network Conditions}
\label{sec:design-network}
\todo{Four slices (Table~\ref{tab:impairments}). Unlike the three-tier
design of~\cite{mosaic2026}, where each tier paired an impairment
with a matched encoder configuration, slices B--D here hold the
encoder fixed at \gls{ldp} \gof{}=4 so that loss, latency and their
interaction are separable; the \gls{ldp} \gof{}=16 block provides the
delivery-unit contrast. Latency and loss values anchored to 3GPP 5G service
performance targets~\cite{3gpp22261}.}

\begin{table}[t]
  \caption{Impairment slices. \todo{fill values}}
  \label{tab:impairments}
  \centering
  \small
  \begin{tabular}{@{}lllll@{}}
    \toprule
    Slice & Loss & Added latency & Encoder & Role \\
    \midrule
    A & 0\%        & \SI{0}{\milli\second}      & \glsentryshort{ra}, \gof{}=32  & Clean reference ($\qc$) \\
    B & \todo{x}\% & \SI{0}{\milli\second}      & LDP, \gof{}=4  & Loss only \\
    C & 0\%        & \todo{y}\,\si{\milli\second} & LDP, \gof{}=4 & Latency only \\
    D & \todo{x}\% & \todo{y}\,\si{\milli\second} & LDP, \gof{}=4 & Interaction \\
    \bottomrule
  \end{tabular}
\end{table}

\subsection{Subjective Experiment}
\label{sec:method}
\paragraph{Protocol.}
\todo{ITU-T P.910 \gls{acr}~\cite{p910}, five-point scale, rated
attribute: smoothness of playback. Each participant views all three
sequences under all four slices; session $\approx$\SI{30}{\minute};
calibration phase on Slice~A; presentation order counterbalanced
(addressing the limitation noted in~\cite{mosaic2026}).}

\paragraph{Participants and screening.}
\todo{$N=50$; demographics; visual-acuity and \gls{xr}-familiarity
screening; outlier rejection per ITU-R BT.500~\cite{bt500} /
ITU-T P.913~\cite{p913}.}

\paragraph{Ethics.}
\todo{Approval via Plataforma Brasil / CEP, protocol number
\anon[withheld for review]{XXXXX}; informed consent; ACM policy on
research involving human participants.}

\subsection{Data Processing and Statistical Analysis}
\label{sec:design-analysis}
\todo{Per-condition \mos with bootstrap 95\% \glspl{ci}; \gls{clmm}
with participant and content as random effects and a latency$\times$loss
interaction term; \gof contrast; leave-one-out evaluation of the pooled
$\alpha$; \gls{plcc}, \gls{srocc}, \gls{rmse} on held-out conditions.
Per-\gof telemetry (deadline outcome, completion time) aligned with
each rating session for the Layer~1 analysis.}
